\chapter{Getting Started}
\label{ch:getting-started}

This chapter builds a small VNET web jail backed by ZFS. The examples assume a
pool named \command{zroot}, an external interface already configured by the
host, and a FreeBSD 15.1 host.

\section{Create a project}

Make a working directory and generate a starter manifest:

\begin{lstlisting}[language=bash]
mkdir -p ~/jails/web-stack
cd ~/jails/web-stack
blackship armada init
\end{lstlisting}

Replace the generated content with a minimal configuration:

\begin{lstlisting}[language=TOML]
[config]
data_dir = "/var/blackship"
zfs_enabled = true
zpool = "zroot"
dataset = "blackship"

[[jails]]
name = "web"
release = "15.1-RELEASE"
hostname = "web.internal"

[jails.network]
vnet = true
networks = ["default"]
ip = "10.0.1.10"

[jails.healthcheck]
enabled = true

[[jails.healthcheck.checks]]
name = "nginx"
command = "service nginx status"
target = "jail"
interval = 30
timeout = 10
retries = 3
\end{lstlisting}

\section{Bootstrap the userland}

\begin{lstlisting}[language=bash]
blackship bootstrap 15.1-RELEASE
blackship releases list
\end{lstlisting}

On FreeBSD 16 and newer, official distribution sets are no longer the normal
source and Blackship uses fingerprint-verified pkgbase automatically.
\command{--pkgbase} opts into that path on supported earlier branches.

\section{Create the named network}

The default epair backend uses an \command{if\_bridge}:

\begin{lstlisting}[language=bash]
blackship network create default \
  --subnet 10.0.1.0/24 \
  --gateway 10.0.1.1 \
  --bridge blackship0
\end{lstlisting}

To use netgraph instead, choose a distinct bridge name:

\begin{lstlisting}[language=bash]
blackship network create default \
  --subnet 10.0.1.0/24 \
  --gateway 10.0.1.1 \
  --backend netgraph \
  --bridge ngbridge0
\end{lstlisting}

Do not create both examples under the same network name. Chapter
\ref{ch:networking-storage} explains the backend tradeoffs and host routing.

\section{Validate and preview}

\begin{lstlisting}[language=bash]
blackship check
blackship up web --dry-run
blackship --format json up web --dry-run
\end{lstlisting}

The plan should identify the full jail name, root path, address, hook counts,
current state, and whether a ZFS dataset will be created. Resolve every
validation error before continuing.

\section{Start and inspect the jail}

\begin{lstlisting}[language=bash]
blackship up web
blackship ps
blackship exec web -- uname -a
blackship console web
\end{lstlisting}

Install and enable a service inside the jail:

\begin{lstlisting}[language=bash]
blackship exec web -- pkg install -y nginx
blackship exec web -- sysrc nginx_enable=YES
blackship exec web -- service nginx start
blackship health web
\end{lstlisting}

For a reproducible service image, put those installation steps in a Jailfile
instead of treating the running jail as the build record; see Chapter
\ref{ch:images}.

\section{Publish a port}

PF anchors isolate Blackship-managed redirects from the rest of the host's PF
configuration:

\begin{lstlisting}[language=bash]
blackship expose web --port 8080 --internal 80
blackship ports web
\end{lstlisting}

The host PF configuration must load Blackship's anchors. A typical
\file{/etc/pf.conf} contains:

\begin{lstlisting}
rdr-anchor "blackship"
anchor "blackship"
\end{lstlisting}

Run \command{blackship setup} after a PF restart, or at boot through the
provided rc integration, to reapply persisted networks and forwards.

\section{Stop and remove resources}

\begin{lstlisting}[language=bash]
blackship down web --dry-run
blackship down web
\end{lstlisting}

\command{down} stops the running jail and releases runtime resources while
preserving its root for future starts. Use \command{rm --volumes} only when the
jail and its associated storage should be removed.

\begin{shipwarning}[Destructive cleanup]
\command{snapshot rollback --force}, \command{rm --volumes}, forced garbage
collection, release deletion, and recursive ZFS destruction can discard data.
Take a snapshot or export first, verify the resolved full jail and dataset names,
and prefer Blackship commands over manual host deletion.
\end{shipwarning}

\section{Next steps}

Add a database service with \setting{depends\_on = ["database"]}, move package
installation into a Jailfile, configure a health-driven QoS transition, and run
\command{blackship supervise}. Those features are developed in the following
chapters.
